\section{Schedule and Model Details}
\label{sec:appendix_schedule}

\paragraph{Schedule model verification.} Substituting $n = 1$ into Eq.~\ref{eq:production_schedule} yields $t_{1,I} = t_0 + (1/k)\ln(2e^{k/\dot{n}_{\max}} - 1) \approx t_0$, confirming that the first unit is produced near the ramp-up midpoint, with $S(t_0) = 0.50$. The piecewise enforcement ($\dot{n} = 0$ for $t < t_c$) affects only the exponentially small tail before $t_c$; $N(t_c)$ is negligible (${\sim}10^{-3}$ units). The counting convention starts from $N(t_0) = 0$, not from facility groundbreaking.

\begin{table}[htbp]
\centering
\caption{Pathway-specific delivery schedules. The gap column shows the additional time ISRU requires relative to Earth delivery of the same unit number.}
\label{tab:production_schedule}
\small
\begin{tabular}{@{}rrrrr@{}}
\toprule
Unit $n$ & $t_{n,E}$ (yr) & $t_{n,I}$ (yr) & $S(t_{n,I})$ & Gap (yr) \\
\midrule
1      & 0.002 & 5.00  & 0.50 & $+$5.00 \\
10     & 0.02  & 5.04  & 0.52 & $+$5.02 \\
100    & 0.20  & 5.34  & 0.66 & $+$5.14 \\
500    & 1.00  & 6.31  & 0.93 & $+$5.31 \\
1{,}000  & 2.00  & 7.34  & 0.99 & $+$5.34 \\
5{,}000  & 10.00 & 15.35 & 1.00 & $+$5.35 \\
10{,}000 & 20.00 & 25.35 & 1.00 & $+$5.35 \\
\bottomrule
\end{tabular}
\end{table}


\paragraph{Deterministic vs.\ MC capital estimates.} The deterministic baseline uses $K = \$50$B (the modal estimate for a first-of-kind vertically integrated ISRU facility), while the MC median is $K = \$65$B. This difference arises because the log-normal distribution is right-skewed: the mode of $\mathrm{LN}(\mu_{\ln}, \sigma_{\ln})$ is $e^{\mu_{\ln} - \sigma_{\ln}^2} = \$46$B, while the median is $e^{\mu_{\ln}} = \$65$B. The deterministic \$50B lies between these values; the MC correctly propagates the full distribution.

